% ===========================================================================
% Armado desde el título de cada nota de docs/decisiones/.
\section{Decisiones de diseño}
\label{anx:decisiones}

Cada nota de \repo{docs/decisiones/} tiene el contexto, las opciones consideradas, la decisión y
sus consecuencias. Una decisión no se edita para cambiarla: se escribe una nueva que la reemplaza.

\begin{longtable}{@{}r p{13.4cm}@{}}
  \toprule
  \textbf{N.\textdegree} & \textbf{Decisión} \\
  \midrule
  \endhead
  001 & Codificación y comportamiento de la instrucción HALT \\
  002 & Formato de trama del protocolo de debug \\
  003 & Baudrate de la UART: 115200 \\
  004 & LOAD como único mecanismo de carga y reprogramación \\
  005 & Cómo se determina la ``memoria de datos usada'' \\
  006 & Debug Unit con dos FSM independientes \\
  007 & Banco de registros: implementación y conflicto de lectura/escritura entre ID y WB \\
  008 & Interfaz de PC: TUI con Textual \\
  009 & Integración de las memorias de lectura sincrónica (BRAM) en el pipeline \\
  010 & Convenciones de código RTL \\
  011 & Detección de load-use conservadora y sin forwarding hacia MEM \\
  012 & Resolución de saltos en EX, sin predicción \\
  013 & Tamaño y mapa de direcciones de las memorias \\
  014 & Implementación de las memorias con BRAM y escritura de la memoria de programa por el puerto B \\
  015 & Accesos desalineados y fuera de rango: se truncan y dejan un aviso \\
  016 & Acceso por byte y media palabra con \emph{byte write enable} de la BRAM \\
  017 & Ensamblador propio, formatos de salida y validación \\
  018 & Ejecución paso a paso sin intervenir el clock \\
  019 & Programas de prueba: estado inicial, variantes con nop y convenciones de sw/ \\
  020 & Modelo de referencia: semántica de la ISA más pipeline ciclo a ciclo \\
  021 & Banco de pruebas incremental: carga, comparación y cortes por etapa \\
  022 & Memoria de programa: sin reset del registro de salida y contenido inicial sin archivo \\
  023 & Interfaz de debug del core y estructura de la etapa IF \\
  024 & Instrucciones no soportadas: decodificación solo por opcode \\
  025 & Filtro de palabras fuera de isa.md en el toolchain \\
  026 & Control de la ALU en ID, sin módulo aparte \\
  027 & Memoria de datos: qué lógica de MEM va en \texttt{dmem} \\
  028 & Regresión del pipeline por lotes en Vivado \\
  029 & Modo del modelo de referencia durante las integraciones de M8 \\
  030 & Corridas sin \texttt{+FINAL} en la regresión por lotes \\
  031 & Mutantes del RTL para la regresión \\
  032 & Testbenches unitarios en la regresión por lotes \\
  033 & Varias corridas en una simulación de \texttt{pipeline\_tb} \\
  034 & Desensamblador: \texttt{rvasm} envuelve el del modelo de referencia \\
  035 & UART de la Debug Unit: interfaz de bytes y FIFO de recepción \\
  036 & FSM de protocolo de la Debug Unit: alcance, puertos y reglas de validación \\
  037 & FSM de ejecución de la Debug Unit: STEP, enable, prioridades y gancho de la carga \\
  038 & Carga de programa en la Debug Unit: contrato con la FSM de ejecución y con el mapa de usadas \\
  039 & Volcado de estado de la Debug Unit: latches, mapa de usadas y fuente del payload \\
  040 & Interfaz de usuario: capas \texttt{model} y \texttt{session}, hilos y corte con la I-53 \\
  041 & Reset del sistema en la placa: botón sincronizado y reset de encendido \\
  042 & Reloj del sistema a 65 MHz con el Clock Wizard para la prueba en placa \\
  043 & Vista del pipeline: instrucción de cada etapa, riesgos inferidos, historial y F4 \\
  044 & Comparador de igualdad para decidir los saltos \\
  045 & Reloj del sistema a 92 MHz, la máxima frecuencia del barrido \\
  \bottomrule
\end{longtable}
